\section*{Figures}

\begin{figure}[htbp]
  \centering
  \IfFileExists{../results/figures/fig1_concept_topology.pdf}{%
    \includegraphics[width=0.95\linewidth]{fig1_concept_topology.pdf}%
  }{\fbox{\parbox{0.9\linewidth}{\centering\vspace{2em}Figure 1 placeholder: \texttt{fig1\_concept\_topology.pdf}\vspace{2em}}}}
  \caption{\textbf{Drug-induced cell-state cycling, detected as a
  topological loop and measured by persistent homology.}
  (a) Untreated cancer cells occupy a uniform transcriptional state;
  under EGFR-TKI pressure they diversify into multiple states with
  reversible transitions between them (cells coloured by state).
  (b) In scRNA-seq, this reversible cycling appears as a closed loop
  in cell-state space (one dot = one cell).
  (c) Persistent homology measures the loop. Around each cell we grow
  a ball; as the ball radius (the filtration scale $\varepsilon$)
  increases, the loop is born when balls touch their neighbours and
  dies when balls fill the loop's interior. The lifespan of this
  loop --- the lifetime of the longest $H_1$ bar --- is the
  \emph{maximum $H_1$ persistence} statistic used throughout this
  work. The mug-and-donut topological-equivalence analogy that
  motivates the use of homology is described in the Background.
  Cross-cohort fold-changes are in \cref{fig:master}.}
  \label{fig:concept}
\end{figure}

\begin{figure}[htbp]
  \centering
  \IfFileExists{../results/figures/fig1_study_overview.pdf}{%
    \includegraphics[width=0.95\linewidth]{fig1_study_overview.pdf}%
  }{\fbox{\parbox{0.9\linewidth}{\centering\vspace{2em}Figure 2 placeholder: \texttt{fig1\_study\_overview.pdf}\vspace{2em}}}}
  \caption{\textbf{Study design and cell-line baseline analysis.}
  (a) UMAP of PC9 erlotinib time course colored by timepoint.
  (b) Leiden clusters.
  (c) EMT score.
  (d) Cell cycle phase.
  (e) Diffusion pseudotime.
  (f) PAGA graph.
  $n = 2{,}942$ cells across six timepoints.}
  \label{fig:study_overview}
\end{figure}

\begin{figure}[htbp]
  \centering
  \begin{subfigure}[b]{0.95\linewidth}
    \centering
    \IfFileExists{../results/figures/fig2_a_persistence_barcodes.pdf}{%
      \includegraphics[width=\linewidth]{fig2_a_persistence_barcodes.pdf}%
    }{\fbox{\parbox{0.9\linewidth}{\centering\vspace{2em}Panel (a) placeholder\vspace{2em}}}}
    \caption{Persistence barcodes by timepoint (top: $H_0$; bottom: $H_1$).}
  \end{subfigure}
  \par\vspace{0.5em}
  \begin{subfigure}[b]{0.75\linewidth}
    \centering
    \IfFileExists{../results/figures/fig2_d_wasserstein_heatmap.pdf}{%
      \includegraphics[width=\linewidth]{fig2_d_wasserstein_heatmap.pdf}%
    }{\fbox{\parbox{0.9\linewidth}{\centering\vspace{2em}Panel (b) placeholder\vspace{2em}}}}
    \caption{Pairwise Wasserstein distances between per-timepoint $H_1$
    persistence diagrams.}
  \end{subfigure}
  \par\vspace{0.5em}
  \begin{subfigure}[b]{0.48\linewidth}
    \centering
    \IfFileExists{../results/figures/fig4_a_barcodes_with_cc.pdf}{%
      \includegraphics[width=\linewidth]{fig4_a_barcodes_with_cc.pdf}%
    }{\fbox{\parbox{0.9\linewidth}{\centering\vspace{3em}Panel (c) placeholder\vspace{3em}}}}
    \caption{Barcodes with cell-cycle genes.}
  \end{subfigure}\hfill
  \begin{subfigure}[b]{0.48\linewidth}
    \centering
    \IfFileExists{../results/figures/fig4_b_barcodes_without_cc.pdf}{%
      \includegraphics[width=\linewidth]{fig4_b_barcodes_without_cc.pdf}%
    }{\fbox{\parbox{0.9\linewidth}{\centering\vspace{3em}Panel (d) placeholder\vspace{3em}}}}
    \caption{Barcodes without cell-cycle genes.}
  \end{subfigure}
  \caption{\textbf{Persistent homology of the PC9 erlotinib discovery
  cohort and Tirosh cell-cycle control.}
  (a) Per-timepoint $H_0$ (top) and $H_1$ (bottom) persistence barcodes
  for the GSE134839 discovery cohort; permutation tests' closest
  approaches to significance are D9 ($p = 0.098$) and D2
  ($p = 0.126$).
  (b) Pairwise Wasserstein distance heatmap between per-timepoint
  $H_1$ persistence diagrams: D0 is separated from all drug-treated
  timepoints in persistence-diagram space.
  (c, d) $H_1$ barcodes recomputed with and without 39 Tirosh
  cell-cycle genes ablated from the HVG set; the persistence-ratio
  $\geq 0.71$ at every timepoint indicates loops are not cell-cycle
  artifacts.}
  \label{fig:discovery_ph}
\end{figure}

\begin{figure}[htbp]
  \centering
  \begin{subfigure}[b]{0.85\linewidth}
    \centering
    \IfFileExists{../results/figures/fig6_d_h1_over_time.pdf}{%
      \includegraphics[width=\linewidth]{fig6_d_h1_over_time.pdf}%
    }{\fbox{\parbox{0.9\linewidth}{\centering\vspace{2em}Panel (a) placeholder\vspace{2em}}}}
    \caption{Max $H_1$ vs.\ days for the PC9 erlotinib and osimertinib
    time courses; osimertinib shows a $2.7\times$ monotonic increase
    from D0 to D14.}
  \end{subfigure}
  \par\vspace{0.5em}
  \begin{subfigure}[b]{0.95\linewidth}
    \centering
    \IfFileExists{../results/figures/fig7_a_pdx_barcodes.pdf}{%
      \includegraphics[width=\linewidth]{fig7_a_pdx_barcodes.pdf}%
    }{\fbox{\parbox{0.9\linewidth}{\centering\vspace{2em}Panel (b) placeholder\vspace{2em}}}}
    \caption{Persistence barcodes for the YU-003 and YU-006 PDX models
    under untreated and osimertinib-residual conditions; both PDX
    models increase $H_1$ persistence under residual disease.}
  \end{subfigure}
  \caption{\textbf{Drug-validation cohorts: osimertinib in PC9 and in
  EGFR-mutant patient-derived xenografts.}
  (a) PC9 max-$H_1$ trajectory under erlotinib vs.\ osimertinib
  (GSE150949). (b) PDX YU-003 and YU-006 untreated vs.\ residual
  barcodes (GSE243562). Together these are the within-system
  monotonic-increase replications behind the cross-cohort summary in
  \cref{fig:master}.}
  \label{fig:validation}
\end{figure}

\begin{figure}[htbp]
  \centering
  \IfFileExists{../results/figures/fig8_master_h1_comparison.pdf}{%
    \includegraphics[width=0.95\linewidth]{fig8_master_h1_comparison.pdf}%
  }{\fbox{\parbox{0.9\linewidth}{\centering\vspace{2em}Master comparison placeholder\vspace{2em}}}}
  \caption{\textbf{Master comparison across all biological systems.}
  (a) Max $H_1$ bar chart across all conditions. (b) Max $H_1$ vs.\
  days for all drug-treated systems, with patient-naive baseline as a
  reference line. Drug treatment pushes cell line and PDX topology
  toward patient-tumor-like complexity.}
  \label{fig:master}
\end{figure}

\begin{figure}[htbp]
  \centering
  \IfFileExists{../results/figures/figS_harmony_batch_correction.pdf}{%
    \includegraphics[width=\textwidth]{figS_harmony_batch_correction.pdf}%
  }{\fbox{\parbox{0.9\linewidth}{\centering\vspace{2em}Harmony figure placeholder\vspace{2em}}}}
  \caption{\textbf{Harmony batch correction sensitivity.}
  (a) Max $H_1$ vs.~days of osimertinib exposure in the PC9 cell line cohort (GSE150949), computed on raw PCA (blue) and on Harmony-corrected embeddings (pink, dashed). Both show monotonic growth from D0 to D14.
  (b) Max $H_1$ per PDX group (untreated vs.~residual disease) in GSE243562, raw (blue) vs.~Harmony-corrected (pink). The YU-006 pre/post effect is preserved after batch correction.}
  \label{fig:harmony}
\end{figure}

\begin{figure}[htbp]
  \centering
  \begin{subfigure}[b]{0.95\linewidth}
    \centering
    \IfFileExists{../results/figures/figS_maynard_barcodes.pdf}{%
      \includegraphics[width=\linewidth]{figS_maynard_barcodes.pdf}%
    }{\fbox{\parbox{0.9\linewidth}{\centering\vspace{2em}Panel (a) placeholder\vspace{2em}}}}
    \caption{Maynard 2020 EGFR-mutant longitudinal cohort
    (14~patients, 13{,}244 cells): persistence barcodes ($H_0$ top,
    $H_1$ bottom) pooled across treatment-naive (TN), persister (PER),
    and progressive disease (PD) stages. Pooled max $H_1$:
    $5.37 \to 7.13 \to 8.05$; three of three matched pre/post biopsies
    increase post-treatment (median $1.60\times$).}
  \end{subfigure}
  \par\vspace{0.5em}
  \begin{subfigure}[b]{0.65\linewidth}
    \centering
    \IfFileExists{../results/figures/fig7_d_kim_atlas_barcodes.pdf}{%
      \includegraphics[width=\linewidth]{fig7_d_kim_atlas_barcodes.pdf}%
    }{\fbox{\parbox{0.9\linewidth}{\centering\vspace{2em}Panel (b) placeholder\vspace{2em}}}}
    \caption{Treatment-naive lung adenocarcinoma reference from the
    Kim et al.~2020 atlas (GSE131907, $n = 1{,}500$ epithelial cells
    subsampled from 5{,}000 tLung cells): persistence barcodes
    establishing the patient-tumor max $H_1 = 4.77$ reference used in
    \cref{fig:master}.}
  \end{subfigure}
  \caption{\textbf{Patient-cohort validation across two independent
  scRNA-seq studies.}
  (a) Maynard 2020 14-patient EGFR-mutant longitudinal cohort:
  max~$H_1$ rises monotonically TN $\to$ PER $\to$ PD with
  3/3 matched pre/post pairs increasing.
  (b) Kim 2020 LUAD atlas: persistence barcodes for the
  treatment-naive primary-tumor epithelium reference. See
  Supplementary Table~S11 for full per-patient Maynard
  trajectories.}
  \label{fig:patient}
\end{figure}

\clearpage
\section*{Tables}

\begin{table}[htbp]
  \centering
  \caption{\textbf{Per-timepoint $H_1$ summary (GSE134839, erlotinib).}
  Permutation $p$-values from 500 within-cell gene-label shuffles.}
  \label{tab:per_timepoint}
  \begin{tabularx}{0.85\linewidth}{l r r r r}
    \toprule
    Timepoint & $n$ cells & Max $H_1$ & \#\,$H_1$ features & Perm.\ $p$ (500) \\
    \midrule
    D0  & 1560 & 1.47 & 1393 & 0.394 \\
    D1  &  324 & 1.32 &  255 & 0.552 \\
    D2  &  230 & 1.77 &  124 & 0.126 \\
    D4  &  200 & 1.37 &   54 & 0.144 \\
    D9  &  379 & 1.77 &  204 & 0.098 \\
    D11 &  249 & 1.08 &  124 & 0.266 \\
    \bottomrule
  \end{tabularx}
\end{table}

\begin{table}[htbp]
  \centering
  \caption{\textbf{Validation cohort $H_1$ permutation tests
  (100 permutations, 1{,}000-cell subsamples).} Values from 1{,}500-cell
  subsample used in main analysis; permutation tests used 1{,}000-cell
  subsamples (see Methods). Per-cohort max $H_1$ values vary across
  subsample sizes (1{,}000 vs.\ 1{,}500 cells) due to the strong
  subsample dependence of persistent homology on cell count
  (Supplementary Table~S7, CV 8.9--23.5\%); values reported in
  Table~\ref{tab:validation} use the 1{,}000-cell permutation-test
  subsample for internal consistency with the $p$-value column. Note
  that PC9 osimertinib D0 max $H_1$ is reported as 1.41 in the main
  analysis (5{,}000-cell subsample) and 1.74 here (1{,}000-cell
  permutation subsample); these are different subsamples of the same
  underlying data and are not inconsistent.}
  \label{tab:validation}
  \begin{tabularx}{0.85\linewidth}{l l r l}
    \toprule
    Dataset & Group & Observed $H_1$ & $p$-value \\
    \midrule
    PC9 + Osimertinib & D0        & 1.74 & 0.030 \\
    PC9 + Osimertinib & D7        & 3.08 & $< 0.01$ \\
    PC9 + Osimertinib & D14       & 3.78 & $< 0.01$ \\
    PDX YU-006        & Untreated & 3.62 & $< 0.01$ \\
    PDX YU-006        & Residual  & 6.08 & $< 0.01$ \\
    PDX YU-003        & Residual  & 3.85 & $< 0.01$ \\
    \bottomrule
  \end{tabularx}
\end{table}

\begin{table}[htbp]
  \centering
  \caption{\textbf{Cell-cycle ablation controls.}
  Top: Tirosh cell-cycle gene removal from HVG set (39 of 97 Tirosh
  genes present; full discovery-cohort timepoints from GSE134839).
  Ratio is max $H_1$ after / before ablation; ratios $\geq 0.71$
  indicate loops are not cell-cycle artifacts.
  Bottom: stringent \texttt{sc.pp.regress\_out} of S-score and G2M-score
  on 1{,}000-cell subsamples across five cohorts. Ratio is max $H_1$
  after regression / original; every ratio $> 1$ indicates the $H_1$
  signal is preserved after removing the dominant cell-cycle variance
  axis.}
  \label{tab:cc_controls}
  \begin{tabularx}{0.95\linewidth}{l r r r l}
    \toprule
    Group & $n$ cells & $H_1$ before & $H_1$ after & Ratio \\
    \midrule
    \multicolumn{5}{l}{\emph{Tirosh cell-cycle gene ablation (GSE134839 erlotinib)}} \\
    \midrule
    D0  & 1560 & 1.47 & 1.35 & 0.92 \\
    D1  &  324 & 1.32 & 1.28 & 0.97 \\
    D2  &  230 & 1.77 & 1.57 & 0.89 \\
    D4  &  200 & 1.37 & 0.97 & 0.71 \\
    D9  &  379 & 1.77 & 1.39 & 0.78 \\
    D11 &  249 & 1.08 & 1.20 & 1.11 \\
    \midrule
    \multicolumn{5}{l}{\emph{Stringent S/G2M regression (1{,}000-cell subsamples)}} \\
    \midrule
    Erl D0               & 1000 & 1.27 & 2.85 & 2.25 \\
    Erl D9               &  379 & 1.77 & 1.89 & 1.07 \\
    Erl D11              &  249 & 1.08 & 4.37 & 4.05 \\
    PDX YU-006 untreated & 1000 & 2.89 & 3.67 & 1.27 \\
    PDX YU-006 residual  & 1000 & 4.65 & 5.96 & 1.28 \\
    \bottomrule
  \end{tabularx}
\end{table}
